\chapter{Discussion}
In this chapter, we reflect on the experimental results presented in Chapter 11. We analyze the specific failure modes of the extraction and classification modules, discuss the fundamental trade-offs between system accuracy and computational cost, and identify the limitations and threats to the validity of our findings, following established methodologies for empirical software and NLP evaluation \cite{buttcher2016information, manning2008introduction}.

\section{Error Analysis}
While the system demonstrated robust performance across both extraction and retrieval tasks, a qualitative analysis of the errors provides insights into the inherent challenges of processing semi-structured job data \cite{sarawagi2008information, tamburri2020data}.

\subsection{Extraction Failures and Implicit Information}
The lower recall for lexical skill extraction highlights a significant challenge: \textit{implicit skill requirements}. Many job postings omit foundational skills that are assumed for a role. This "semantic gap" is a known issue in recruitment informatics, where requirements are often latent rather than explicit \cite{giabelli2021skills, qin2018enhancing}.

In the case of date extraction, errors primarily occurred in "rolling" recruitment cycles. Our current regex-based Layer A often failed to normalize these into a precise temporal range, a limitation common in rule-based Information Extraction systems when faced with high linguistic variability \cite{grishman2015information, appelt1999introduction}.

\subsection{Ambiguity in Semantic Mapping}
The gap between Top-1 and Top-3 accuracy for title normalization reveals the impact of professional ambiguity. Generic titles lack sufficient semantic density for the bi-encoder to map them to a specific ESCO occupation without analyzing the full text body. This challenge of entity resolution and disambiguation is central to modern taxonomy alignment research \cite{zhang2022survey, sayfullina2018learning}.

\section{Trade-offs: Accuracy vs Cost}
The design of our tiered architecture was motivated by the need to balance performance with operational feasibility, a core principle in "Green AI" and efficient deep learning \cite{schwartz2020green, menghani2023efficient}.

\subsection{The Cost of Perfection}
While Large Language Models (LLMs) offer superior zero-shot capabilities, the marginal gain in accuracy comes at a significant cost in both latency and financial expenditure \cite{hsieh2023distilling, bender2021dangers}. By delegating the majority of the workload to local, small-scale models, we achieved a "good enough" accuracy level suitable for production environments while minimizing reliance on expensive external APIs \cite{menghani2023efficient, sanh2019distilbert}.

\subsection{Latency and User Experience}
The search latency of 45ms is critical for a responsive user interface. A system relying on real-time LLM inference for search-time semantic expansion would likely exceed the thresholds for interactive systems \cite{buttcher2016information}. Our use of pre-computed HNSW indices and RRF-based fusion allows the system to deliver semantic results with the efficiency of traditional lexical search \cite{malkov2018efficient, cormack2009reciprocal}.

\section{Threats to Validity}
Despite the positive results, several factors may affect the generalizability and robustness of this study, as is common in Information Retrieval and machine learning research \cite{manning2008introduction, thakur2021beir}.

\subsection{Internal Validity: Annotation Bias}
The evaluation relies on a "Gold Standard" annotated by experts. While inter-annotator agreement was high, subjectivity in labeling "relevant" skills remains a threat, a known challenge in creating benchmarks for specialized domains \cite{jurafsky2023speech, thakur2021beir}.

\subsection{External Validity: Generalizability}
Our dataset is primarily focused on the Tech and Business sectors within the North American and European contexts, driven by the ESCO framework \cite{esco2017handbook}. The effectiveness of the system might degrade when applied to different linguistic domains or regions with different occupational structures \cite{kurekova2015using}.

\subsection{Construct Validity: Metric Selection}
We used standard IR metrics like MAP and NDCG. However, these metrics do not always correlate perfectly with real-world user satisfaction in job discovery, where factors like "freshness" and company reputation may play significant roles \cite{baeza2011modern, lacic2020recommender}.
